\chapter{Declaración de uso de IA generativa}\label{anx:DeclaracionIA}
% Tutoría anotada el 2026-11-09: declaración breve del uso de IA generativa en el
% desarrollo del producto y en la redacción de la memoria (herramientas, para qué
% se usaron y cómo se revisó lo generado). Formato según la normativa de TFG de la
% ESI (pendiente de consultar).
% BORRADOR del agente (2026-11-24): el alumno lo revisa, lo corrige y lo hace suyo;
% las frases sobre lo que hace el alumno solo son válidas si se corresponden con lo
% que realmente ha hecho.
\pendiente{adaptar al formato que exija la normativa de TFG de la ESI}

El uso de agentes de IA es el objeto de este trabajo. Se declara a continuación en qué se han empleado:

\begin{itemize}
	\item \textbf{Herramienta.} Claude Code, de Anthropic, en su aplicación de escritorio y en la extensión para VS Code, con GitHub Spec Kit. Los \emph{commits} elaborados con su ayuda lo indican en su mensaje, junto con el modelo empleado.
	\item \textbf{Producto.} La IA ha propuesto y redactado la constitución, las especificaciones, los planes y las tareas, y ha generado el código y las pruebas. El alumno ha tomado las decisiones, ha aprobado cada artefacto y ha revisado cada hito.
	\item \textbf{Memoria.} La IA ha elaborado borradores de los capítulos a partir de los documentos del proyecto y de fuentes citadas. \pendiente{el alumno describe aquí su revisión y reescritura, una vez hecha}
	\item \textbf{Perfiles de usuario.} Las proto-personas y sus respuestas se han elaborado con IA. Son supuestas y no se presentan como datos.
	\item \textbf{Control de lo generado.} Las normas se contrastan con su publicación oficial, las afirmaciones técnicas llevan fuente, el código se verifica con pruebas y un pipeline, y las decisiones del alumno constan en el registro del proceso (anexo~\ref{anx:Registro}).
\end{itemize}

La responsabilidad sobre el contenido de la memoria y sobre el producto corresponde al autor.
